%%
%% This is file `sample-sigconf.tex',
%% generated with the docstrip utility.
%%
%% The original source files were:
%%
%% samples.dtx  (with options: `all,proceedings,bibtex,sigconf')
%% 
%% IMPORTANT NOTICE:
%% 
%% For the copyright see the source file.
%% 
%% Any modified versions of this file must be renamed
%% with new filenames distinct from sample-sigconf.tex.
%% 
%% For distribution of the original source see the terms
%% for copying and modification in the file samples.dtx.
%% 
%% This generated file may be distributed as long as the
%% original source files, as listed above, are part of the
%% same distribution. (The sources need not necessarily be
%% in the same archive or directory.)
%%
%%
%% Commands for TeXCount
%TC:macro \cite [option:text,text]
%TC:macro \citep [option:text,text]
%TC:macro \citet [option:text,text]
%TC:envir table 0 1
%TC:envir table* 0 1
%TC:envir tabular [ignore] word
%TC:envir displaymath 0 word
%TC:envir math 0 word
%TC:envir comment 0 0
%%
%% The first command in your LaTeX source must be the \documentclass
%% command.
%%
%% For submission and review of your manuscript please change the
%% command to \documentclass[manuscript, screen, review]{acmart}.
%%
%% When submitting camera ready or to TAPS, please change the command
%% to \documentclass[sigconf]{acmart} or whichever template is required
%% for your publication.
%%
%%
\documentclass[sigconf]{acmart}
%%
%% \BibTeX command to typeset BibTeX logo in the docs
\AtBeginDocument{%
  \providecommand\BibTeX{{%
    Bib\TeX}}}

%% Rights management information.  This information is sent to you
%% when you complete the rights form.  These commands have SAMPLE
%% values in them; it is your responsibility as an author to replace
%% the commands and values with those provided to you when you
%% complete the rights form.
\setcopyright{acmlicensed}
\copyrightyear{2018}
\acmYear{2018}
\acmDOI{XXXXXXX.XXXXXXX}
%% These commands are for a PROCEEDINGS abstract or paper.
\acmConference[Conference acronym 'XX]{Make sure to enter the correct
  conference title from your rights confirmation email}{June 03--05,
  2018}{Woodstock, NY}
%%
%%  Uncomment \acmBooktitle if the title of the proceedings is different
%%  from ``Proceedings of ...''!
%%
%%\acmBooktitle{Woodstock '18: ACM Symposium on Neural Gaze Detection,
%%  June 03--05, 2018, Woodstock, NY}
\acmISBN{978-1-4503-XXXX-X/2018/06}


%%
%% Submission ID.
%% Use this when submitting an article to a sponsored event. You'll
%% receive a unique submission ID from the organizers
%% of the event, and this ID should be used as the parameter to this command.
%%\acmSubmissionID{123-A56-BU3}

%%
%% For managing citations, it is recommended to use bibliography
%% files in BibTeX format.
%%
%% You can then either use BibTeX with the ACM-Reference-Format style,
%% or BibLaTeX with the acmnumeric or acmauthoryear sytles, that include
%% support for advanced citation of software artefact from the
%% biblatex-software package, also separately available on CTAN.
%%
%% Look at the sample-*-biblatex.tex files for templates showcasing
%% the biblatex styles.
%%

%%
%% The majority of ACM publications use numbered citations and
%% references.  The command \citestyle{authoryear} switches to the
%% "author year" style.
%%
%% If you are preparing content for an event
%% sponsored by ACM SIGGRAPH, you must use the "author year" style of
%% citations and references.
%% Uncommenting
%% the next command will enable that style.
%%\citestyle{acmauthoryear}


%%
%% end of the preamble, start of the body of the document source.
\begin{document}

%%
%% The "title" command has an optional parameter,
%% allowing the author to define a "short title" to be used in page headers.
\title{Dependability analysis of a Java open-source project}

%%
%% The "author" command and its associated commands are used to define
%% the authors and their affiliations.
%% Of note is the shared affiliation of the first two authors, and the
%% "authornote" and "authornotemark" commands
%% used to denote shared contribution to the research.

\author{Matthias Gerhard}
\affiliation{%
 \institution{Università degli Studi di Salerno}
 \city{Salerno}
 \state{Campania}
 \country{Italy}}

\author{Max Timur Sönmez}
\affiliation{%
  \institution{Università degli Studi di Salerno}
  \city{Salerno}
  \state{Campania}
  \country{Italy}}

%%
%% By default, the full list of authors will be used in the page
%% headers. Often, this list is too long, and will overlap
%% other information printed in the page headers. This command allows
%% the author to define a more concise list
%% of authors' names for this purpose.
\renewcommand{\shortauthors}{Trovato et al.}

%%
%% The abstract is a short summary of the work to be presented in the
%% article.
\begin{abstract}
  A clear and well-documented \LaTeX\ document is presented as an
  article formatted for publication by ACM in a conference proceedings
  or journal publication. Based on the ``acmart'' document class, this
  article presents and explains many of the common variations, as well
  as many of the formatting elements an author may use in the
  preparation of the documentation of their work.
\end{abstract}

%%
%% The code below is generated by the tool at http://dl.acm.org/ccs.cfm.
%% Please copy and paste the code instead of the example below.
%%
\begin{CCSXML}
<ccs2012>
   <concept>
       <concept_id>10002978.10003022.10003023</concept_id>
       <concept_desc>Security and privacy~Software security engineering</concept_desc>
       <concept_significance>500</concept_significance>
       </concept>
   <concept>
       <concept_id>10011007.10011074.10011099.10011102.10011103</concept_id>
       <concept_desc>Software and its engineering~Software testing and debugging</concept_desc>
       <concept_significance>500</concept_significance>
       </concept>
   <concept>
       <concept_id>10003456.10003457.10003490.10003503.10003505</concept_id>
       <concept_desc>Social and professional topics~Software maintenance</concept_desc>
       <concept_significance>300</concept_significance>
       </concept>
   <concept>
       <concept_id>10011007.10011074.10011099.10011692</concept_id>
       <concept_desc>Software and its engineering~Formal software verification</concept_desc>
       <concept_significance>100</concept_significance>
       </concept>
 </ccs2012>
\end{CCSXML}

\ccsdesc[500]{Security and privacy~Software security engineering}
\ccsdesc[500]{Software and its engineering~Software testing and debugging}
\ccsdesc[300]{Social and professional topics~Software maintenance}
\ccsdesc[100]{Software and its engineering~Formal software verification}

%%
%% Keywords. The author(s) should pick words that accurately describe
%% the work being presented. Separate the keywords with commas.
\keywords{Do, Not, Use, This, Code, Put, the, Correct, Terms, for,
  Your, Paper}

\received{20 February 2007}
\received[revised]{12 March 2009}
\received[accepted]{5 June 2009}

%%
%% This command processes the author and affiliation and title
%% information and builds the first part of the formatted document.
\maketitle

\section{Introduction}

Modern software development increasingly relies on complex microservice architectures, often utilizing robust frameworks such as Spring Boot and event-streaming platforms like Apache Kafka. While these technologies accelerate functional development, they often introduce complexity that makes rigorous quality assurance and formal verification difficult. Open-source repositories frequently demonstrate functional correctness but lack the software dependability layers—such as mutation testing, strict static analysis, and formal specifications—required for production-grade systems.

This paper presents a case study on retrofitting dependability measures onto an existing open-source artifact. We selected the \texttt{ecommerce-api-spring-kafka} repository \cite{ecommerceRepo}, a representative microservices application implementing an e-commerce backend. While the original repository provided a functional implementation of shopping cart and order processing logic, it lacked a comprehensive automated testing strategy and formal behavioral specifications.

Our work follows a systematic approach to software hardening, addressing the following key areas:

\begin{enumerate}
    \item \textbf{Environmental Stabilization:} We addressed immediate build reproducibility issues by upgrading the Docker base image to \texttt{eclipse-temurin:17-jdk} and pinning Kafka/Zookeeper versions to resolve startup failures in the containerized environment.
    \item \textbf{Automated Quality Gates:} We established a continuous integration (CI) pipeline using GitHub Actions to enforce build verification on every push.
    \item \textbf{Static and Dynamic Analysis:} We integrated a suite of static analysis tools (SpotBugs, PMD, Checkstyle) and dynamic coverage analysis (JaCoCo) to establish a quality baseline.
    \item \textbf{Advanced Verification:} We extended the test suite based on coverage data and applied Mutation Testing (PiTest) to assess the robustness of our tests. Additionally, we applied Java Modeling Language (JML) annotations to core services to formalize system behavior, identifying architectural constraints imposed by framework-heavy designs (e.g., Lombok) that complicate formal verification.
\end{enumerate}

The remainder of this paper details the methodology used to integrate these tools, presents the empirical metrics gathered (coverage improvements, mutation scores), and discusses the challenges of applying formal specifications to modern Spring Boot applications.

% --- 2. DEPENDABILITY GOALS ---
\section{Dependability Goals}

The primary objective of this study was to transform a functional prototype into a robust, maintainable, and verifiable software artifact. We defined four specific dependability goals to guide our interventions:

\begin{itemize}
    \item \textbf{Reproducibility \& Stability:} The build environment must be deterministic. The original repository suffered from runtime failures due to unpinned dependencies and incompatible base images. Our goal was to ensure the application builds and launches consistently across different environments.
    \item \textbf{Code Quality Baseline:} We aimed to establish a quantitative baseline for code quality using industry-standard static analysis tools (SpotBugs, PMD, Checkstyle) to identify potential bugs, bad practices, and style violations before applying fixes.
    \item \textbf{Test Completeness \& Robustness:} Merely having tests is insufficient. We aimed to increase meaningful code coverage (excluding generated code) and verify the quality of the test suite itself using Mutation Testing. The goal was to ensure that tests are sensitive enough to detect real regressions.
    \item \textbf{Behavioral Specification:} We sought to formalize the business logic of core services using the Java Modeling Language (JML). The goal was to document strict preconditions and postconditions for critical methods (e.g., \texttt{ShoppingCartService}), creating a clear contract for future development, even if full formal verification is limited by framework constraints.
\end{itemize}

% --- 3. METHODOLOGY ---
\section{Methodology}

Our methodology followed a phased approach, starting with environmental stabilization and progressing toward advanced verification techniques.

\subsection{Phase 1: Environment Stabilization \& CI}
Before addressing the code, we stabilized the execution environment.
\begin{itemize}
    \item \textbf{Containerization:} We updated the Docker base image from \texttt{openjdk:17-jdk-slim} to \texttt{eclipse-temurin:17-jdk} to resolve compatibility issues with the project's build environment.
    \item \textbf{Dependency Pinning:} The original \texttt{docker-compose.yml} used \texttt{latest} tags for Kafka and Zookeeper, causing startup failures. We pinned these services to version \texttt{3.3.1} to ensure stability.
    \item \textbf{Continuous Integration:} We implemented a GitHub Actions workflow that triggers on every push to the \texttt{develop} branch. This pipeline executes \texttt{mvn clean verify} to enforce that the application compiles and passes all tests automatically.
\end{itemize}

\subsection{Phase 2: Static Analysis Integration}
We integrated three distinct static analysis tools into the Maven build lifecycle to detect defects early.
\begin{itemize}
    \item \textbf{Tool Selection:} We employed \textbf{SpotBugs} for bug patterns, \textbf{PMD} for programming flaws, and \textbf{Checkstyle} for adherence to coding standards.
    \item \textbf{Execution:} These tools were configured to run via \texttt{mvn verify}, generating XML reports (\texttt{checkstyle-result.xml}, \texttt{pmd.xml}, etc.).
    \item \textbf{Strategy:} We categorized findings by severity (Critical, Medium, Low) but intentionally applied no fixes at this stage. This allowed us to preserve the original state as a baseline for measuring the impact of later testing and mutation efforts.
\end{itemize}

\subsection{Phase 3: Coverage \& Mutation Testing}
To improve the reliability of the test suite, we used a data-driven approach.
\begin{itemize}
    \item \textbf{Coverage Analysis (JaCoCo):} We integrated JaCoCo to measure instruction and branch coverage. Crucially, we configured exclusions for generated code (e.g., \texttt{*MapperImpl.class}) to prevent false negatives in our metrics. The initial analysis revealed a test coverage of 48.4\% with only 14.4\% of branches tested.
    \item \textbf{Test Expansion:} Based on the JaCoCo reports, we extended the unit test suite, specifically targeting uncovered branches in the domain logic.
    \item \textbf{Mutation Testing (PiTest):} We implemented PiTest to evaluate the quality of our tests. We encountered "unkillable mutants" in Lombok-generated code (e.g., Builders, Getters/Setters) and AOP aspects. Consequently, we refined our configuration to exclude boilerplate code, focusing mutation analysis on the core domain logic where semantic meaning exists.
\end{itemize}

\subsection{Phase 4: Formal Specification (JML)}
In the final phase, we applied JML annotations to document the behavior of critical services.
\begin{itemize}
    \item \textbf{Scope:} We targeted the \texttt{ShoppingCartService} and \texttt{OrderService}, defining invariants, preconditions (e.g., non-null arguments), and postconditions (e.g., total price calculations).
    \item \textbf{Constraints:} We found that the project's heavy reliance on Spring Boot and Lombok prevented the use of OpenJML's static verification (\texttt{esc}) and runtime checking. The "magic" bytecode generation of these frameworks obscures the plain Java logic required by JML tools. Therefore, JML served primarily as a formal documentation layer rather than an automated verification gate.
\end{itemize}

% --- 4. EMPIRICAL RESULTS ---
\section{Empirical Results}

This section presents the quantitative data gathered from our dependability interventions, focusing on code coverage, mutation scores, and the feasibility of formal specifications in the target environment.

\subsection{Static Analysis Baseline}
The integration of SpotBugs, PMD, and Checkstyle provided an initial quality baseline. The analysis was executed via the Maven build cycle, generating XML reports for all modules.
\begin{itemize}
    \item \textbf{Findings:} The tools successfully categorized issues by severity (Critical, Medium, Low).
    \item \textbf{Outcome:} These reports served as a static baseline. No automated fixes were applied at this stage to ensure that subsequent dynamic testing metrics reflected the state of the original logic rather than auto-corrected code.
\end{itemize}

\subsection{Test Coverage (JaCoCo)}
We utilized JaCoCo to measure the effectiveness of the unit test suite. To ensure the metrics reflected relevant business logic, we configured the tool to exclude generated artifacts, specifically `*MapperImpl.class` (MapStruct implementations) and `**/generated/**` directories.

\textbf{Coverage Metrics:}
\begin{itemize}
    \item \textbf{Line Coverage:} 48.4\%
    \item \textbf{Branch Coverage:} 14.4\%
\end{itemize}

These results highlighted a significant gap in the testing of conditional logic (branch coverage). The analysis revealed that while the "happy path" was often tested, edge cases and error handling branches in the service layer remained largely unverified.

\subsection{Mutation Testing (PiTest)}
To assess the quality of the test suite beyond simple line execution, we applied PiTest. The initial runs identified numerous "unkillable" mutants arising from boilerplate code (e.g., Lombok builders, getters/setters). After filtering these non-semantic mutations, we focused on the domain packages.

Table \ref{tab:mutation} summarizes the mutation coverage improvements achieved after extending the test suite based on the initial reports.

\begin{table}[h]
  \centering % Optional: centers the table on the page
  \begin{tabular}{l c p{4cm}}
    \toprule
    \textbf{Package} & \textbf{Mut. Score} \\
    \midrule
    domain.address & 25\% $\rightarrow$ 100\% \\
    domain.shoppingcart & 60\% $\rightarrow$ 77\% \\
    infrastructure.aspect & 50\% \\
    infrastructure.exceptions & 60\% \\
    \bottomrule
  \end{tabular}
  \caption{Mutation Testing Results by Package}
  \label{tab:mutation}
\end{table}

The improvement in `domain.address` (reaching 100\%) validates that value objects, often overlooked, are easy to perfect once identified. However, the `infrastructure` packages demonstrate the limitation of mutation testing when applied to framework-heavy code (Aspects/Exceptions), where mutations often result in invalid bytecode or behavior identical to the original.

\subsection{Formal Specification (JML)}
We successfully annotated the core service methods with JML to document invariants and contracts.
\begin{itemize}
    \item \textbf{Specification Scope:} Annotations were applied to \emph{ShoppingCartService} (method: \emph{buildCartItems}, \emph{updateTotalPrice}) and \emph{OrderService} (method: \emph{getShippingAddress}, \emph{buildOrder}).
    \item \textbf{Verification Result:} \textit{N/A (Qualitative Only)}.
\end{itemize}

While the specifications improve code readability and contract clarity, empirical verification (via \texttt{openjml --esc}) was not possible. The project's architecture—relying heavily on compile-time bytecode generation (Lombok) and runtime dependency injection (Spring)—creates an Abstract Syntax Tree (AST) that standard JML tools cannot parse without significant refactoring. Thus, the result for JML is a successful \textit{specification} effort, but a failed \textit{verification} attempt due to toolchain incompatibility.
% --- 5. DISCUSSION & INSIGHTS ---
\section{Discussion \& Developer Insights}

The process of retrofitting dependability measures onto the \emph{ecommerce-api-spring-kafka} repository revealed distinct challenges and trade-offs inherent in modern Java development.

\subsection{The Framework vs. Verification Conflict}
A major insight from this study is the friction between "developer productivity" tools and "formal verification" tools.
\begin{itemize}
    \item \textbf{Lombok \& JML Incompatibility:} The project utilizes Project Lombok to reduce boilerplate (getters, setters, builders). While this improves code readability, it obscures the source code from JML tools, which expect plain Java syntax. We found that formal verification (using OpenJML) is effectively impossible in a heavy-Lombok environment without delombok-ing the code first—a step that complicates the build pipeline significantly.
    \item \textbf{Implication:} Teams must decide early: prioritize rapid prototyping (Lombok) or verifiable correctness (JML). Retrofitting verification later is cost-prohibitive.
\end{itemize}

\subsection{Mutation Testing: Signal vs. Noise}
Our PiTest integration demonstrated that mutation testing is valuable but requires aggressive tuning.
\begin{itemize}
    \item \textbf{The "Noise" Problem:} Initially, framework code (AOP aspects, Exception classes) generated hundreds of unkillable mutants. For example, mutating an Exception class constructor rarely changes program behavior, leading to false negatives in the report.
    \item \textbf{The Solution:} We found that restricting mutation testing strictly to the \texttt{domain} packages provided high-value insights (improving coverage from 60\% to 77\%), whereas running it on `infrastructure` layers was computationally expensive and analytically low-yield.
\end{itemize}

\subsection{Infrastructure as a Prerequisite}
We observed that dependability is not just about code. The initial inability to run the application due to unpinned Docker images and Kafka versions proved that \textit{build reproducibility} is the first step of reliability. A test suite is useless if the CI pipeline cannot spin up the environment deterministically.

% --- 6. CONCLUSION ---
\section{Conclusion}

This paper presented a case study on hardening an existing open-source microservices application. By establishing a stabilized Docker environment and a GitHub Actions CI pipeline, we created a reproducible foundation for quality assurance.

Our empirical results show that while standard unit testing provided a starting point, it left significant gaps in branch coverage (only 14.4\% covered). The introduction of Mutation Testing was critical in exposing these gaps, driving improvements in the domain layer where semantic logic resides. However, our attempt to introduce formal verification via JML highlighted a significant gap in the ecosystem: modern frameworks like Spring Boot and Lombok are currently resistant to standard formal verification tools due to their reliance on runtime injection and bytecode generation.

We conclude that for existing "modern" Java applications, a combination of \textbf{Containerization}, \textbf{Strict Static Analysis}, and \textbf{Targeted Mutation Testing} offers the highest return on investment for dependability, while formal verification remains a specialized technique best suited for core algorithmic modules rather than full-stack web architectures.

% --- 7. REFERENCES ---
\bibliographystyle{ACM-Reference-Format}
\bibliography{references} 

\end{document}
\endinput
%%
%% End of file `sample-sigconf.tex'.
